\begin{figure}[tb]
\centering
\begin{tikzpicture}[x=1.65cm,y=1.2cm,>=Stealth,font=\small]
 \fill[blue!8] (2,2)--(3.7,3.7)--(3.7,2)--cycle;
 \draw[->] (0,0)--(4.25,0) node[right] {$\lambda$};
 \draw[->] (0,0)--(0,4.05) node[above] {eigenvalue};
 \draw[gray,dashed,thick] (0.35,0.35)--(3.85,3.85);
 \draw[blue!60!black,very thick] (0.35,0.35)--(2,2)--(3.9,2);
 \draw[dotted] (2,0)--(2,2);
 \draw[dotted] (0,2)--(2,2);
 \node[below] at (2,0) {$2-p^2$};
 \node[left] at (0,2) {$2-p^2$};
 \draw (3,0.045)--(3,-0.045) node[below] {$2$};
 \node[gray,rotate=36,above] at (1.1,1.1) {original eigenvalue};
 \node[blue!60!black,below] at (3.15,1.9) {$\min(\lambda,2-p^2)$};
 \draw[<->,thick] (3.5,2.06)--(3.5,3.45);
 \node[right,align=left] at (3.57,2.8)
   {observation\\precision};
 \draw[<->] (2,-0.65)--(3.95,-0.65);
 \node[below,align=center] at (2.98,-0.65)
   {active eigendirections\\$\lambda>2-p^2$};
\end{tikzpicture}
\caption{A spectral cap, shown near the upper edge. The posterior interaction
$W_p$ retains eigenvalues below $2-p^2$ and flattens those above it.
The removed amount $(\lambda-(2-p^2))_+$ is the precision of the Gaussian
observation in that eigendirection. The drawing is schematic; its geometry
illustrates the defining spectral functions, not an eigenvalue density.}
\label{fig:spectral-cap}
\end{figure}
